\documentclass[twocolumn,amsmath,amssymb]{snp}
\pagestyle{empty}
\usepackage{graphicx}
\usepackage{bm}
\topmargin 1.5 cm
\textwidth14.5cm
\textheight20cm
\oddsidemargin0.7cm
\columnsep0.2in

\begin{document}

\title{{\Large Dynamical Origin of Monopole-Quadrupole Mixing in Highly Deformed $^{24}$Mg}}
\author{\large Abhishek}
\email{abhishek@ph.iitr.ac.in}
\affiliation{Department of Physics, Indian Institute of Technology Roorkee,
Roorkee - 247667, INDIA}
\maketitle

\section*{Introduction}
Quadrupole deformation couples the isoscalar giant monopole resonance (ISGMR)
to the $K=0$ branch of the isoscalar giant quadrupole resonance, producing a
low-energy monopole component beside the main compression mode
\cite{kvasil2016}. High-resolution $\alpha$ scattering from strongly prolate
$^{24}$Mg found pronounced strength near 13--14 MeV and a broader structure at
higher energy \cite{bahini}. Wavelet analysis separated a 10--18 MeV
monopole-quadrupole-coupling (MQC) region from the 18--24 MeV main ISGMR region
\cite{bahini}.

Quasiparticle Random-Phase Approximation (QRPA) identifies the mixed modes through simultaneous IS0 and $K=0$ IS2 strength~\cite{bahini}; here we complement this eigenmode picture by calculating the reciprocal off-diagonal $E0\leftrightarrow E2(K=0)$ response in Time-Dependent Hartree-Fock (TDHF).

\begin{figure}[t]
\centering
\includegraphics[width=\columnwidth]{../figures/fig1_e0_experiment_two_edf.pdf}
\caption{\label{fig:experiment}Measured $^{24}$Mg IS0 strength and TDHF with
SLy5 and SkM*. The ordinate is in fm$^4$ MeV$^{-1}$. Theory uses the
$\sum_i r_i^2$ convention without a fitted scale or energy shift. Both EDFs
separate the MQC and main-ISGMR regions but miss the measured fragmentation.
Here $\Gamma_{\rm sm}=1$ MeV.}
\end{figure}

\begin{figure*}[t]
\centering
\includegraphics[width=0.74\textwidth]{../figures/fig2_combined_mqc_two_edf.pdf}
\caption{\label{fig:dynamics}TDHF responses: SLy5 diagonal strengths $S_{00}$
(a) and $S_{20}$ (b), reciprocal signed mixed responses $R_{20,00}$ and
$R_{00,20}$ (c), and $R_{20,00}$ for both EDFs (d). All ordinates are in
fm$^4$ MeV$^{-1}$, but the $F_{00}$ and $F_{20}$ normalizations differ, so
panel magnitudes are not directly comparable. Shading marks 9--18 MeV;
$\Gamma_{\rm sm}=0.2$ MeV.}
\end{figure*}

\section*{Method}
The axial prolate SLy5 and SkM* minima have $Q_{20}=36.532$ and
35.568 fm$^2$, corresponding to $\beta_2=0.319$ and 0.311 in the
Sky3D total-mass convention. After aligning the symmetry axis with $z$,
weak boosts excite the isoscalar monopole and $K=0$ quadrupole operators,
$F_{00}=r^2/\sqrt{4\pi}$ and $F_{20}=\sqrt{5}r^2Y_{20}$.
Calculations use Sky3D v1.2 \cite{abhishek2024} with the SLy5 and SkM*
parametrizations \cite{forces}, no pairing, and a nonperiodic $24^3$
grid with 1-fm spacing.

Following the boost $\exp(-i\eta F_A)$, the TDHF state is evolved to
$T=18000$ fm/$c$ with $\eta=5\times10^{-5}$, while both $F_{00}(t)$
and $F_{20}(t)$ are monitored. The diagonal and mixed responses are
obtained from
\begin{equation}
 \begin{split}
 \chi^{(-)}_{BA}(E)&=\frac{1}{\eta\hbar}\int_0^T\!dt\,
 e^{-iEt/\hbar}w(t)\,\delta F_B(t),\\
 R_{BA}(E)&=\pi^{-1}\operatorname{Im}\chi^{(-)}_{BA}(E),
 \end{split}
\end{equation}
where $A,B\in\{00,20\}$ and
$w(t)=\exp[-\Gamma_{\rm sm}t/(2\hbar)]$.
For $A=B$, $R_{AA}$ is the usual strength function, while $A\ne B$
gives the signed off-diagonal response. Linearity and finite-box
stability were checked with reduced-amplitude and larger-box calculations.

\section*{Results and discussion}

Figure~\ref{fig:dynamics} shows that the low-energy monopole structure is
accompanied by a strong $K=0$ quadrupole response. For both SLy5 and SkM*,
the diagonal $E0$ and $E2(K=0)$ strengths and the off-diagonal response peak
at nearly the same energy. Thus, a monopole perturbation in deformed
$^{24}$Mg directly drives the quadrupole degree of freedom associated with
the MQC mode. The independently calculated $E0\!\rightarrow\!E2$ and
$E2\!\rightarrow\!E0$ responses satisfy Onsager--Kubo reciprocity
\cite{kubo1957} within numerical accuracy and remain unchanged for a reduced
boost strength, confirming a linear cross-channel response.

The two EDFs give the same qualitative separation between the low-energy
mixed mode and the higher-energy main ISGMR, although their detailed spectra
differ. Existing box-size tests show that centroids and energy-window
integrals are stable, whereas individual narrow TDHF lines are not.

As shown in Fig.~\ref{fig:experiment}, the calculated centroids of the
10--18 MeV MQC and 18--25 MeV main-ISGMR regions lie within about $0.6$ MeV
of experiment without an energy shift or fitted normalization. TDHF,
however, overconcentrates the low-energy strength and does not reproduce the
observed fragmentation, indicating missing correlations beyond mean field.

\section*{Conclusion}
TDHF time evolution shows that an $E0$ excitation of deformed $^{24}$Mg
directly induces the $K=0$ quadrupole motion associated with the
monopole--quadrupole mixed mode. The effect is robust for both SLy5 and
SkM* and is consistent with the experimentally identified low-energy MQC
region. The detailed fragmentation, however, remains beyond mean-field TDHF.

{\footnotesize
\begin{thebibliography}{5}
\bibitem{kvasil2016} J. Kvasil \textit{et al.}, Phys. Rev. C \textbf{94},
064302 (2016).
\bibitem{bahini} A. Bahini \textit{et al.}, Phys. Rev. C \textbf{105},
024311 (2022); \textbf{113}, 064313 (2026).
\bibitem{abhishek2024} Abhishek \textit{et al.}, Comput. Phys. Commun.
\textbf{301}, 109239 (2024).
\bibitem{forces} E. Chabanat \textit{et al.}, Nucl. Phys. A \textbf{635}, 231
(1998); J. Bartel \textit{et al.}, Nucl. Phys. A \textbf{386}, 79 (1982).
\bibitem{kubo1957} R. Kubo, J. Phys. Soc. Jpn. \textbf{12}, 570 (1957).
\end{thebibliography}
}

\end{document}
